% Quarantine-only candidate statement.
% Date: 29 August 2026.
% This file is not a live-book input and has not been Rethlas-verified.
% It makes no novelty or priority claim.

\begin{lemma}
\label{lem:p213-variation-graph-branch-separation}
Let $K$ be a complete non-Archimedean field, possibly trivially valued,
let $Y$ be a good separated normal boundaryless $K$-analytic space that is
purely one-dimensional, and let $y\in Y^{[2,3]}$.  Fix a finite analytic
moment block $F_0$ near $y$ and a finite family $A$ of branches of $Y$ at
$y$.

Inside any prescribed open neighborhood of $y$, one can first choose an
open tree $W_*$ about $y$ and append a finite block $S$ of analytic units,
and then perform the boundary-avoiding source localization entirely inside
$W_*$.  This produces a compact analytic neighborhood $V\subset W_*$ of
$y$ and a moment
\[
       F=(F_{\mathrm{bd}},F_0,S):V\longrightarrow
       \mathbf G_{m,K}^{N}
\]
with the following properties.  If
\begin{equation}
 \Omega:=\mathbf R^N\setminus
 F_{\mathrm{trop}}\bigl(\partial(V/Y)\bigr),
 \qquad
 U:=F_{\mathrm{trop}}^{-1}(\Omega),
 \label{eq:p213-variation-local-omega}
\end{equation}
then $U$ is an open neighborhood of $y$ in $Y$ and
$F_{\mathrm{trop}}|_U:U\to\Omega$ is proper.  Near
$v:=F_{\mathrm{trop}}(y)$, the tropical image is a finite rational fan.
Every branch in $A$ on which at least one coordinate of $F_0$ varies is
mapped to a ray distinct from the ray of every other source branch.

More precisely, let $b\in A$ be such a branch.  Every sufficiently short
relatively compact open subsegment $\rho_b$ of its target ray that does not
contain $v$ admits an ambient open neighborhood $O_b\subset\Omega$ such
that
\begin{equation}
 O_b\cap F_{\mathrm{trop}}(U)=\rho_b
 \label{eq:p213-variation-ambient-tube}
\end{equation}
and $F_{\mathrm{trop}}^{-1}(O_b)$ is contained in a section of $b$, apart
from source subtrees on which every coordinate of $F$ is locally constant.
In particular, $F_{\mathrm{trop}}^{-1}(O_b)$ is a genuine open source
tube.  The exact inverse image of the ray itself is not asserted to be
open.
\end{lemma}

\begin{lemma}
\label{lem:p213-euclidean-one-dimensional-current}
Let $\Omega\subset\mathbf R^N$ be open, with its standard integral affine
structure, and let $\Sigma\subset\Omega$ be a locally finite embedded
rational polyhedral complex of pure dimension one.  Let $R$ be a positive
$\mathrm d'$- and $\mathrm d''$-closed supercurrent of bidimension $(1,1)$
on $\Omega$ whose support is contained in $\Sigma$.  After subdividing
$\Sigma$, there are constants $b_e\in\mathbf R_{\geq0}$, one for every
open edge $e$, such that
\begin{equation}
                         R=\sum_e b_e[e].
 \label{eq:p213-euclidean-edge-decomposition}
\end{equation}
The sum is locally finite and $[e]$ is the integration current calibrated
by the primitive integral tangent of $e$.  At every vertex $v\in\Omega$
one has
\begin{equation}
                    \sum_{e\ni v}b_eu_{e,v}=0,
 \label{eq:p213-euclidean-vertex-balance}
\end{equation}
where $u_{e,v}$ is the primitive outgoing tangent.  In particular, $R$
has no additional current supported at a vertex.
\end{lemma}

\begin{proposition}
\label{prop:p213-normal-curve-arbitrary-valuation-current}
Let $K$ be a complete non-Archimedean field, possibly trivially valued,
and let $Y$ be a connected good separated normal boundaryless
$K$-analytic space that is purely one-dimensional.  Every positive
$\mathrm d'$- and $\mathrm d''$-closed ordinary Chambert--Loir--Ducros
current $T$ of bidegree $(0,0)$ on $Y$ is a nonnegative constant multiple
of the integration current:
\begin{equation}
                         T=c\,\delta_Y,
             \qquad c\in\mathbf R_{\geq0}.
 \label{eq:p213-normal-curve-arbitrary-classification}
\end{equation}
Consequently $T$ is a Bedford--Taylor current in the sense of
\cite[Definition~8.6.2]{CLD25}.
\end{proposition}

\begin{proposition}
\label{prop:p213-curve-nonsmooth-locus-dichotomy}
Let $K$ be a complete non-Archimedean field and let $Y$ be an irreducible
good strict reduced $K$-analytic space that is purely one-dimensional.
Exactly one of the following alternatives occurs.
\begin{enumerate}
\item The quasi-smooth locus of $Y$ is nonempty.  Its complement is a
closed locally finite set of rigid points.  If $Y$ is boundaryless, the
same assertion holds with ``rig-smooth'' in place of ``quasi-smooth''.
\item The quasi-smooth locus is empty.
\end{enumerate}
Thus, under these hypotheses, the condition that the non-rig-smooth locus
is not contained in a locally finite set of rigid points is equivalent to
the normalized irreducible curve being nowhere rig-smooth.  In
particular, there is no intermediate case consisting of isolated
nonrigid failures of rig-smoothness.
\end{proposition}

\begin{example}
\label{ex:p213-nowhere-smooth-normal-curve}
Let $K=\mathbf F_p((t))$, let $L=K(t^{1/p})$, and regard
$Y=(\mathbf A^1_L)^{\mathrm{an}}$ as a $K$-analytic curve.  Then $Y$ is
good, separated, strict, boundaryless, irreducible, normal, and purely
one-dimensional, but it is nowhere quasi-smooth over $K$.  Hence the
second alternative in
\cref{prop:p213-curve-nonsmooth-locus-dichotomy} actually occurs, and its
nonsmooth locus contains nonrigid points.
\end{example}

\begin{remark}
\label{rem:p213-variation-graph-quarantine-scope}
The preceding statements are quarantined candidates.  Their proof uses
ordinary analytic moments and ordinary CLD currents.  It does not assert
the analogous classification for weakly smooth forms presented only by
harmonic tropicalizations.  It also does not by itself settle the passage
across conductor points on a reduced nonnormal curve over a trivially
valued field.
\end{remark}
